\documentclass[10pt,a4paper]{amsart}
\usepackage[margin=19mm]{geometry}
\usepackage{amsmath,amssymb,mathtools}
\usepackage[scr=rsfs,cal=euler]{mathalfa}
\usepackage{xfrac}
\usepackage[unicode,hidelinks,bookmarks=false]{hyperref}
\newcommand{\ie}{i.e.\ }
\newcommand{\cf}{cf.\ }
\newcommand{\resp}{respectively\ }
\newcommand{\Subsection}{\subsection}
\providecommand{\nobreakdash}{\nobreak\hbox{-}}
\setlength{\emergencystretch}{2em}
\multlinegap0pt
\usepackage{amssymb}
\usepackage{mathtools}
\usepackage{xcolor}
\newtheorem{lemma}{Lemma}
\title{Forcing monochromatic subdivisions}
\author{Gabriel Collado, Dhruv Mubayi}
\date{Source: arXiv:2609.00636v1 (CC BY 4.0)}
\begin{document}
\maketitle

\noindent\emph{Source and licence.} Forcing monochromatic subdivisions. Version arXiv:2609.00636v1 (CC BY 4.0), distributed by arXiv under the Creative Commons Attribution 4.0 International licence, http://creativecommons.org/licenses/by/4.0/, as recorded in the arXiv licence metadata for this exact version and shown on the abstract page https://arxiv.org/abs/2609.00636v1. Full licence notice: \texttt{licenses/arxiv-2609.00636-CC-BY-4.0.txt}.\par\smallskip

\noindent\textit{Editorial scope.} Counted unit: the article's lemma lem:robust (source0.tex lines 231--241). The article's complete original proof of it is delivered with it (source0.tex lines 243--340, one \texttt{proof} environment of 98 lines). No mathematics is written, repaired or re-derived: the statement and the proof are the article's own lines, in the article's own notation, and no retained line is substituted or re-flowed.  Retained uncounted as the source-local context the proof needs: lines 211--220.  Nothing was omitted from any retained span, so the package carries no omission record.  No retained line contains a citation, so the wrapper carries no bibliography and no bibliography entry.  No retained line contains a figure environment, an image-inclusion command or drawing code, so no image is delivered and the package declares no asset dependency.  Editorial formatting only, in addition: an AMS \texttt{amsart} preamble that restates the article's own declarations and macros used by the retained lines, namely the environments \texttt{lemma} on separate counters, together with the standard packages those lines need.  The retained environments print in this document as Lemma 1 in order of appearance, whereas the article prints the same items with the numbering of its own full document.  The complete native source of the article as distributed in the arXiv e-print for this exact version is bundled unchanged as \texttt{sources/209046/source0.tex}.
\begin{lemma}\label{lem:EML}
Let $G$ be an $N$-vertex $r$-regular graph with $\lambda(G)\le\lambda$.
If $X,Z\subseteq V(G)$ are disjoint, then
\[
\left|e_G(X,Z)-\frac{r}{N}|X||Z|\right|
   \le \lambda\sqrt{|X||Z|},
\]
where $e_G(X,Z)$ is the number of edges with one endpoint in $X$ and the
other in $Z$.
\end{lemma}

\begin{lemma}\label{lem:robust}
Fix $s\in\mathbb{Z}^+$. Let $T\ge 100$ and put
\[
\eta=\frac{1}{10s\sqrt{T}}.
\]
Let $G$ be an
$N$-vertex $r$-regular graph satisfying
$\lambda(G)\le\eta r$. Then every set $S\subseteq V(G)$ with
$|S|\ge N/s$ contains a set $U\subseteq S$ with $|U|>N/(3s)$ such that
$G[U]$ is a $(1/(10s),T)$-expander.
\end{lemma}

\begin{proof}
Set
\[
a=\frac{N}{25sT}
\qquad\text{and}\qquad
b=\frac{N}{10sT}.
\]
We first claim that
\begin{equation}\label{eq:medium-expansion}
|N_{G[S]}(X)|\ge T|X|
\end{equation}
for every $X\subseteq S$ satisfying $a\le |X|\le a+b$.

Suppose otherwise, and set
\[
Y=N_{G[S]}(X),
\qquad
Z=S\setminus(X\cup Y).
\]
There are no edges between $X$ and $Z$. Moreover,
\[
|X|\le a+b=\frac{7N}{50sT}
\qquad\text{and}\qquad
|Y|<T|X|\le\frac{7N}{50s}.
\]
Consequently,
\[
|Z|\ge \frac{N}{s}-\frac{7N}{50sT}-\frac{7N}{50s}
     \ge \frac{N}{3s},
\]
where the last inequality uses $T\ge100$. By
Lemma~\ref{lem:EML} and $e_G(X,Z)=0$,
\[
\frac{r}{N}|X||Z|
 \le \lambda(G)\sqrt{|X||Z|}.
\]
It follows that
\[
|X||Z|
 \le \frac{\lambda(G)^2}{r^2}N^2
 \le \eta^2N^2
 =\frac{N^2}{100s^2T}.
\]
Since $|Z|\ge N/(3s)$, this gives
\[
|X|\le\frac{3N}{100sT}<\frac{N}{25sT}=a,
\]
a contradiction. This proves \eqref{eq:medium-expansion}.

We now perform a pruning procedure. Start with $U_0=S$. As long as there
is a nonempty set $X_i\subseteq U_{i-1}$ such that
\[
|X_i|\le b
\qquad\text{and}\qquad
|N_{G[U_{i-1}]}(X_i)|<T|X_i|,
\]
delete $X_i$ and set $U_i=U_{i-1}\setminus X_i$.

We claim that the procedure deletes fewer than $a$ vertices in total.
Otherwise, let $k$ be the first index for which
\[
W=X_1\cup\cdots\cup X_k
\]
has size at least $a$. By the minimality of $k$,
\[
a\le |W|<a+b.
\]
Furthermore,
\[
N_{G[S]}(W)
 \subseteq \bigcup_{i=1}^k N_{G[U_{i-1}]}(X_i),
\]
and therefore
\[
|N_{G[S]}(W)|
 \le \sum_{i=1}^k |N_{G[U_{i-1}]}(X_i)|
 < T\sum_{i=1}^k|X_i|
 =T|W|.
\]
This contradicts \eqref{eq:medium-expansion}.

Let $U$ be the set remaining when the procedure terminates. Then
\[
|U|>|S|-a
 \ge\frac{N}{s}-\frac{N}{25sT}
 >\frac{N}{3s}.
\]
By the stopping rule,
\[
|N_{G[U]}(X)|\ge T|X|
\]
for every nonempty $X\subseteq U$ with $|X|\le b$; the same inequality is
trivial for $X=\varnothing$. Since
\[
\frac{|U|}{10sT}\le\frac{N}{10sT}=b,
\]
the graph $G[U]$ is a $(1/(10s),T)$-expander.
\end{proof}

\end{document}
